这份参考覆盖 Croqtile 编程的方方面面：

\begin{itemize}[leftmargin=1.6em]
\item \textbf{程序结构}：Croqtile 程序如何组织——\co@__co__@ 函数、主机侧 API 与 C++ 互操作。
\item \textbf{带形状数据}：Croqtile 面向多维张量的类型系统——形状、\co@MDSpan@、\co@spanned data@、符号维度与动态形状。
\item \textbf{循环与并行性}：单程序多数据并行、循环控制、分块与迭代构造。
\item \textbf{数据搬运}：DMA 语句、用 \co@chunkat@ 分块、重排与高级数据搬运模式。
\item \textbf{C++ 嵌入}：输入／输出约定、\co@__cpp__@ 内联块、宏与预处理器集成。
\item \textbf{优化模式}：tileflow 优化、异步 DMA 模式、多缓冲与性能调优技术。
\end{itemize}
